% chapters/ch4-sec5.tex -- 5. Recapitulation and Qualitative Features of the Analytical Results;
% 5.1 Bengen's Maxima; 5.2 Model Rocket Design Optimization; 5.2.1-5.2.6 (PDF 634-645, printed
% pp. 602-613), from the heading "5." (mid-PDF 634) to the end of the chapter text. No numbered
% equations; Figure 16.
\section{Recapitulation and Qualitative Features of the Analytical Results}\label{ch4:sec:5}

Now that the methods of Chapters~\ref{ch2} and~\ref{ch3}, and the present chapter, have placed at
our disposal methods for predicting the dynamic behavior, drag, and altitude capability of any
model rocket we may turn our attention to the qualitative features of the results obtained in
order to ascertain the relative importance of each characteristic, or \emph{parameter}, of the
vehicle to its overall altitude performance. The flight dynamics and aerodynamics chapters each
presented criteria for good design practice in relation to their own subject matter. This section
will attempt to integrate those criteria with some that are basic to the ballistic characteristics
of model rockets and to arrive at some fundamental, unified design rules.

\subsection{Bengen's Maxima}\label{ch4:sec:5.1}

If one were to construct a graph of \emph{total altitude as a function of initial mass for various
values of the drag constant $k$}, given a specific model rocket engine type, a curious effect would
be noticed: there is an optimum, nonzero weight that gives maximum altitude performance. If the
model is either heavier \emph{or lighter} than the optimum weight, its altitude capability will be
less than it would be at the optimum. The existence of such optimum points was first noted by
Bengen in 1965. Malewicki, working independently, published explicit examples of graphs displaying
the altitude maxima in 1966 as model rocket manufacturer's reports. It is also true, as you can see
from the sample plot presented in Figure 19,\ednote{The 1973 text reads ``Figure 19''; the chapter
has no Figure 19. Figure~\ref{ch4:fig:16} is evidently meant: it is the chapter's one plot of
altitude against liftoff mass for several values of $k$, and its caption calls it an example of
Bengen's maxima. Neither the errata nor the supplements correct this; the reference
is kept as printed.} that \emph{the optimum weight
decreases as the value of $k$ decreases} until, in the limit as $k$ approaches zero, the optimum
weight also approaches zero. It \emph{is} actually possible, for most values of $k$ to which
practical model rockets can be built, to build a rocket that is \emph{too light} for maximum
altitude performance! It is a remarkable coincidence, as noted by Stine in his \emph{Handbook of
Model Rocketry} even prior to Bengen and Malewicki, that many model rockets just happen to be
constructed near the ballistically optimum point.

The effect which produces Bengen's maxima can be explained qualitatively in the following manner.
A typical model rocket gains most of its altitude during the coasted portion of the flight. The
coasted altitude increment depends heavily upon two things (other than the deceleration due to
gravity, which is constant for all model rockets): the \emph{burnout momentum} and the \emph{drag}
due to the burnout velocity. It actually depends analytically on the burnout mass, the burnout
velocity, and the drag parameter $k$. If the vehicle is extremely light, its burnout velocity will
be large since the burnout velocity increases monotonically with decreasing weight (i.e., the
burnout velocity always increases as the weight decreases). The low weight, however, will prevent
the burnout \emph{momentum} from being very large. The \emph{drag}, on the other hand, will be
extremely high due to its quadratic dependence on velocity. As a result, the momentum of the rocket
is \emph{dissipated}, or used up, very rapidly and the coasted altitude is disappointingly low.
This is the same effect that insures that a whiffle ball or ping-pong ball cannot travel very far
regardless of how hard it is hit. The extreme limit of this behavior would be the example of a
feather being shot out of a cannon: the muzzle velocity is tremendous but the momentum of the
feather is almost instantly overcome by its large drag.

If we now have a rocket that is very \emph{heavy}, its burnout velocity will be relatively
\emph{low}. Its large mass, however, insures that its burnout \emph{momentum} will be about the
same as that of the lighter rocket (actually, with the exception of the momentum lost to drag and
gravity during the burning phase, the two rockets will have \emph{identical} values of burnout
momentum; this must be true since the total impulse of the engine is just the total momentum
change imparted to the rocket by its engine). The low burnout velocity insures that aerodynamic
drag will dissipate the burnout momentum much more slowly than is the case for the light rocket.
The deceleration of \emph{gravity}, however, will dissipate the burnout momentum at a rate much
\emph{greater} than in the case of the light rocket---and again the altitude performance suffers.

Thus, the optimum liftoff mass lies somewhere between the two extremes. Since the methods of
Section~\ref{ch4:sec:2} of this chapter have given us closed-form, algebraic solutions for
altitude performance in cases of vertical flight, it is theoretically possible to obtain a
closed-form expression for the ballistically optimum liftoff mass. When one attempts this in
practice, however, it is found that the resulting expression is far more complicated than the
original total altitude function and requires a prohibitive amount of time and effort to solve in
any given case. While such an expression satisfies the letter of the definition, ``closed-form
solution'', it literally covers pages when written out; it does not satisfy the \emph{intent} of
the definition as used in this book, which requires a practical closed-form solution to be quicker
and easier to use than numerical or graphical methods from which the same result could be
obtained.

The most convenient method yet found for determining the optimum liftoff mass of a model rocket is
to consult a graph of the kind described at the beginning of this section. Sets of such plots are
currently available from Estes Industries, Incorporated of Penrose, Colorado as Technical Report
TR-10, \emph{Altitude Prediction Charts}; and from the Centuri Engineering Company of Phoenix,
Arizona as Technical Information Report TIR-100, \emph{Model Rocket Altitude Performance}.
Ambitious modelers can prepare their own charts for engine types not listed in these reports,
using the methods of Section~\ref{ch4:sec:2} except in very high-performance cases, where numerical
interval methods must be used due to drag divergence at transonic speeds. It is also suggested
that calculations be restricted to values of liftoff mass greater than the mass of the loaded
engine alone (which is the least possible mass of any actual model) and that values of $k$ be kept
above a reasonable minimum associated with the lowest possible drag coefficient to which a model
rocket can, as a practical matter, be designed.

\subsection{Model Rocket Design Optimization}\label{ch4:sec:5.2}

With the sum total of the information presented throughout this volume ready to our hand, we are
now in a position to formulate a general design procedure which, if adhered to, will enable the
model rocketeer to formulate designs which are optimized from the standpoint of the altitude they
are capable of achieving with the specific model rocket engine for which they are designed to be
flown. It is important to remember that our criterion of optimization is \emph{altitude
capability} alone. The material presented in this book, properly applied, permits the synthesis of
rocket designs that may be considered ``superior'', or even ``best'', for applications depending
solely upon maximum altitude achieved. Examples of such applications are micrometeorological
sounding and the National Association of Rocketry's altitude and payload competition events. There
exist applications in the field of model rocketry, however, for which the rocket with the greatest
altitude capability is not necessarily the best one. Examples that come immediately to mind in this
category are the NAR parachute, boost/glider, and rocket glider duration events, in which the
model's altitude capability must be balanced against the need for a parachute or aerodynamic
surfaces large enough to insure a slow descent and a good probability of the descending model's
being caught by thermals. With the caution that our optimization techniques are valid only for
models in which altitude capability is the sole criterion of excellence, then, we present the
following design optimization procedure.

\subsubsection{Initial Design Definition}\label{ch4:sec:5.2.1}

Given the purpose for which your model is to be used and the model rocket engine with which it is
to be flown, make a rough sketch of your first impression of what it should look like. Refine the
sketch as necessary, adding dimensions (the dimensions are, of course, approximate and may be
changed later as necessary). Simplicity is the rule at this stage in the design: avoid \emph{all}
unnecessary complication of shape and arrangement. A way-out, Buck Rogers design may make a fine
sport, demonstration, or display model, but it will rarely if ever come out looking very good when
judged solely on the basis of its altitude capability. Also avoid unnecessary \emph{size}, making
the design as compact as possible consistent with adequate provisions for a good, reliable recovery
system and any payload that is to be carried.

\subsubsection{Drag Coefficient}\label{ch4:sec:5.2.2}

Using the techniques presented in Section~\ref{ch3:sec:6} of Chapter~\ref{ch3}, estimate the
coefficient of drag of your model at a zero angle of attack. You should be somewhat conservative
regarding your assumptions of skin-friction drag at this point, both to allow for small
imperfections encountered in finishing the model and to allow for the deterioration of the
rocket's surface that will occur from flight to flight. In designing model rockets for maximum
altitude performance the lowest possible drag coefficient is desirable, and to this end you should
not hesitate to experiment with variations on the initial design formulated in
Section~\ref{ch4:sec:5.2.1}. Avoid the use of launch lugs on high-performance models if at all
possible by launching from a low-friction tower or closed-breech mechanism, or at least try to
incorporate a ``pop-off'' or ``shuck-off'' lug design if a lugless launching device is not
available. Remember that, if you do use a fixed launch lug system, you will probably have to accept
an increment in drag coefficient $(\Delta\CD)_{\mathrm{lug}}$ on the order of 0.25 as given in
Section~\ref{ch3:sec:4.3.3} of Chapter~\ref{ch3}. Since it is sometimes possible to design model
rockets with values of $\CD$ less than 0.4, provided no launch lug is present, you can see that the
presence of a fixed lug exacts a considerable penalty in terms of altitude performance.

\subsubsection{Weight Optimization}\label{ch4:sec:5.2.3}

Use the first-estimate value of $\CD$ obtained in Section~\ref{ch4:sec:5.2.2} to compute the drag
parameter $k$ (it will be recalled that $k$ is equal to $\tfrac{1}{2}\rho\CD A_r$, where $A_r$ is
the reference cross-sectional area and $\rho$ is the density of the atmosphere). If you will be
launching from a site considerably above sea level, or if your model will reach an extremely great
altitude (of which you have a rough idea), be sure to adjust your value of $\rho$ from its
sea-level value of 1.225 kilograms per cubic meter. Remember that $k$, for use with charts such as
Figure~\ref{ch4:fig:16}, must be computed using MKS units. If you have a single-staged,
single-engined model you may then use a Malewicki chart (as plots such as Figure~\ref{ch4:fig:16}
have come to be called) to select the optimum weight for your rocket directly by locating Bengen's
maximum for that particular value of $k$. If you are using clustering or staging, some numerical
experimentation with different values of liftoff mass using the methods described in
Section~\ref{ch4:sec:2} will be required. For some rockets having very low values of $k$, it may
not be possible to achieve the optimum weight (for instance, when the weight of the payload and
engine alone are found to total more than the optimum). In such cases, the model must simply be
built as light as practicable consistent with adequate structural strength. For rockets to which
weight must be added to optimize the liftoff mass, adding nose weight is usually best since this
permits adequate stability to be attained with smaller fins than would otherwise be the case, and
thus lessens drag.

\begin{figure}[htbp]
  \centering
  \includegraphics{figures/v2/ch4/fig16.pdf}
  \caption{Maximum altitude attained in vertical flight vs.\ liftoff mass for models using the Type
    B4 engine; an example of Bengen's maxima. Curves have been plotted for models having several
    values of the drag parameter $k$; the $k$ values, in kg/m, have been noted on the curves. Since
    the line on which all Bengen's maxima lie is to the right of the line representing the initial
    mass of the engine itself (for all values of $k$ to which model rockets can actually be built),
    Bengen's maxima are practical considerations that must be taken into account when designing
    model rockets for maximum altitude performance. Graphs of this type were first compiled by
    Douglas J. Malewicki and are therefore called ``Malewicki charts''.}
  \label{ch4:fig:16}
\end{figure}

\subsubsection{Dynamic Stability Optimization}\label{ch4:sec:5.2.4}

Determine the values of the dynamic parameters $C_1$, $C_2$, $I_L$, and $I_R$ for your tentative
design according to the methods presented in Section~5 of Chapter~\ref{ch2}.\ednote{The 1973
text reads ``Section 5 of Chapter 2''. Section~\ref{ch2:sec:4} of Chapter~\ref{ch2}, Analytical
Determination of the Dynamic Parameters, is evidently meant: it computes $C_1$, $C_2$, $I_L$ and
$I_R$ from the size, shape and mass distribution of a design not yet built, while
Section~\ref{ch2:sec:5} measures them on an actual vehicle. Neither the errata nor the supplements
correct this; the reference is kept as printed.} Adjust
these values by altering the mass distribution, length, fin size, fin shape, and fin placement as
necessary to obtain a static stability margin of about one caliber or slightly more, a damping
ratio $\zeta$ between 0.05 and 0.3, and a desirable value of the natural frequency $\omega_n$. We
would suggest, based on experience to date, that the most desirable natural frequencies for
vehicles of model rocket size lie between $0.002v$ and $0.01v$, where $v$ is the airspeed of the
model given in \emph{centimeters} per second (remember that CGS units are used throughout
Chapter~\ref{ch2}). The equivalent limits for airspeeds given in meters per second are $0.2v$ and
$1.0v$. A lower value usually means the longitudinal moment of inertia is too great for the
corrective moment coefficient, so that the model cannot respond to the moments applied by its fins
with sufficient rapidity for safe and stable flight at average, present-day launch departure
velocities. A higher value usually means the longitudinal moment of inertia is too low, so the
model is easily disturbed, or that the corrective moment coefficient is too large, usually
indicating too-large fins that cause excessive drag. If you wish to use a roll rate to reduce
horizontal dispersion, check to insure that it is not too near the coupled resonant frequency
$\omega_{\mathrm{cres}}$. You may also wish to check the roll rates induced by various assumed
slight fin misalignments to guard against the possibility of resonance due to unintentional roll.

\subsubsection{Reduction of Drag at Angle of Attack}\label{ch4:sec:5.2.5}

Consistent with the requirements of the preceding steps, attempt to select or refine the fin design
such that the increase in drag with angle of attack is minimized. Since, for any given value of
$k$, the reduction in altitude performance due to dynamic oscillations depends strongly on the
ratio of $\epsilon$ to $k$, a model with a relatively small value of this ratio has an advantage
over one with a larger value (provided the two vehicles have identical resistance to in-flight
disturbances). Elliptical fins are often used to help here, since it is known that an
\emph{airplane wing} that is not twisted and has a moderate aspect ratio will have the lowest
induced drag due to lift if its planform is elliptical.

\subsubsection{Philosophy of Design and Flight}\label{ch4:sec:5.2.6}

You have no doubt noticed that the procedures outlined in each of the
Sections~\ref{ch4:sec:5.2.1} through~\ref{ch4:sec:5.2.5} tend to alter the design of the model from
what it was at the previous stage of the procedure. The possibility thus exists that optimization
of the design at any given stage, using information from the previous stage, may to a certain
extent invalidate the figures taken from the previous stage. To ``home in'' on an optimum design to
an arbitrary degree of precision it is therefore generally necessary to perform several design
\emph{iterations}; that is, to go through the process outlined in Sections~\ref{ch4:sec:5.2.1}
through~\ref{ch4:sec:5.2.5} several times, using the status of the design from the previous
iteration, or ``pass'', to perform the next iteration. A great deal of this is done in professional
industry, but the modeler can, at his option, choose the number and precision of his iterations at
will, depending on how precise he desires his optimization to be. For many purposes, of course, a
bare minimum of analysis, not even using all the methods and procedures outlined in this book, will
suffice to produce a reasonably good design. For contest work, however, or record-setting attempts
and scientific payload lofting, many rocketeers will want to take full advantage of the most
advanced and painstaking analytical methods of design available.

No matter how careful the design of his vehicle has been, however, there exist certain conditions
encountered at the launch site which may sometimes force the modeler to depart from strictly
altitude-optimized design to achieve the desired results in actual operation. If a strong
horizontal wind exists at the time of launch, for instance, a model designed for a low-thrust,
long-burning engine will often achieve a higher altitude if launched with a high-thrust,
short-burning engine of identical specific impulse, since it will not then have time to respond
completely to the horizontal wind encounter during the burning phase of flight. An example would be
the substitution of a B14 engine in a model optimized for a B4, in case of strong winds at the
launch site. There also exist instances in which it is wise for the modeler to render his design
deliberately off-optimum in order to remain below cloud cover or within the capabilities of limited
tracking facilities. Such cases, however, are entirely a matter of the rocketeer's intuition and
expertise gained through experience and we do not see fit to discuss them in a book devoted
strictly to the optimization of the actual altitude capability of model rockets.

In closing, the authors of this volume would like to state that it is not their intention to end
all questions of model rocket design with the material presented herein. To the best of our
knowledge and belief based on experience to date, the techniques and information presented in this
book form a valid framework within which each model rocketeer can construct his own philosophy of
altitude optimization. A great deal remains to be said, however, for the value of intuition and
experience, and much experimental and analytical research remains to be done in model rocketry, a
field that is at once a science, a sport, a hobby, and a fine art.
